\chapter{Conclusions}
\label{chap:conclusions}

We solved the ELM LASSO problem
\begin{equation*}
    \min_{\mathbf{w}\in\mathbb{R}^{H}}\;
    f(\mathbf{w}) = \tfrac{1}{2}\,\|\mathbf{X}\mathbf{w}-\mathbf{y}\|_{2}^{2}
    + \lambda_\text{LASSO}\,\|\mathbf{w}\|_{1}
\end{equation*}
with two algorithms whose convergence properties sit at opposite ends
of the nonsmooth-optimisation spectrum. IRLS reformulates the
non-differentiable $\ell_{1}$ penalty as a sequence of smooth quadratic
surrogates and inherits the linear rate of the MM framework; SGPTL
keeps the non-smoothness explicit and pays for it with the
$\Omega(L^{2}/\varepsilon^{2})$ lower bound that limits any first-order
method on convex non-differentiable functions.

\paragraph{What the experiments confirmed.}
On synthetic ELM-style instances the empirical behaviour tracks the
theory of Chapter~\ref{chap:comparison} closely. IRLS produced a
straight line on semilog gap-versus-iteration plots, monotone
non-increasing $f$ at every seed we tried, and a sparse iterate at
convergence with the support recovered to within two percentage points
of the ground truth. SGPTL produced an oscillating
$f(\mathbf{w}_{i})$ stabilised by a monotone record value
$\bar{f}^{i}$, a sublinear gap curve that flattens as the iterations
accumulate, and an iterate with zero exact zeros, in line with the
absence of any explicit thresholding mechanism. The $O(MH^{2})$
pre-computation cost of IRLS was visible on the scalability plot as a
clean cubic slope at $M=5H$, and the work required by each method to reach the same gap on the
convergence instance ($1.5\cdot 10^{-6}$ for IRLS at iteration $100$
versus $\approx 3\cdot 10^{-6}$ for SGPTL at iteration $8000$, an
$80\times$ iteration multiplier for the same accuracy) lines up with
the $O(\log\varepsilon^{-1})$ versus $O(\varepsilon^{-2})$ comparison.
On the moderate problem of Section~\ref{sec:iters-to-eps} IRLS reached
$\varepsilon=10^{-2}$ in $3$~iterations against SGPTL's $664$, and
$10^{-3}$ in $7$ against $8678$; SGPTL reached $10^{-6}$ only near the
end of its $30000$-iteration budget (iteration $29184$), while IRLS
reached it in $101$ iterations. This illustrates how quickly the rate
gap dominates as the accuracy target tightens.

\paragraph{Where the theoretical bounds proved conservative.}
A few observations did not have a clean theoretical counterpart in the
report and are worth recording. The IRLS final gap saturated around
$10^{-10}$ for $\varepsilon_\text{thr}=10^{-12}$, but pushing
$\varepsilon_\text{thr}$ further did not improve the gap further: the
linear solver hits floating-point limits before the safety mechanism
does, so the ``smaller $\varepsilon_\text{thr}$ is always sharper''
intuition has a hard floor. The SGPTL contraction factor $\rho$ proved
highly influential at OLS warm start: the final record gap spreads over
several orders of magnitude across $\rho\in[0.5,0.99]$
(from $2.3\cdot 10^{-8}$ at $\rho=0.7$ to
$1.2\cdot 10^{-2}$ at $\rho=0.99$, see Section~\ref{sec:params}).
Aggressive contraction wins at OLS warm start, with $\rho=0.7$ the
empirical optimum ($2.3\cdot 10^{-8}$), against the $[0.9, 0.95]$ range
the slide recommendation of Section~\ref{sec:dsm_parameters} would
suggest as a generic default. We keep $\rho=0.9$ as the default in the
comparison and scalability runs for consistency with
Section~\ref{sec:convergence}; the qualitative conclusions are
unaffected by the $\rho$ choice. The initial gap estimate $\delta_{0}$
was much less sensitive on the tested OLS-warm-start instance: all
values in $\{0.01,0.05,0.1,0.5,1.0\}f^{*}$ ended in the
$10^{-6}$ range.
Finally, the OLS warm start mandated by
Section~\ref{sec:dsm_algorithm} interacts unfavourably with the
greedy deflection. The current subgradient becomes nearly aligned
with $\mathbf{d}_{i-1}$ within a few tens of iterations,
$\gamma_{i}\to 0$, and the stepsize-restricted
$\beta_{i}=\min(\beta,\gamma_{i})$ multiplies the Polyak step by that
same vanishing factor, so the iterate stops moving before the
travel-distance counter $r$ can ever reach $R$ to trigger a
$\delta$-contraction. We had to add an iteration-count fallback to
the patience mechanism — contract $\delta$ and reset
$\mathbf{d}_{i-1}=\mathbf{0}$ when $\bar{f}^{i}$ has not improved for
$R_\text{iter}=\max(i_\text{max}/100,\,50)$ iterations — to recover
the sublinear progress the theory predicts. The reset is consistent
with the per-step bound~\eqref{eq:proof_perstep}, which holds iterate-by-iterate
without requiring persistent memory.

\paragraph{Recommendation for the ELM LASSO problem.}
On the project's intended regime, $M\gg H$ with $\mathbf{W}_{1}$ fixed
random and the output layer trained offline once, IRLS is the recommended
choice. It has a single numerical knob ($\varepsilon_\text{thr}$),
delivers sparse iterates without post-processing, and reaches
machine-precision accuracy in tens to a few hundred iterations on every
problem size we tested. SGPTL becomes competitive only when $H^{2}\gg
M$, i.e. when the hidden layer is wide and the training set is small
enough that the $O(MH^{2})$ pre-computation cost of IRLS no longer
amortises, or when the dense $H\times H$ matrix
$\mathbf{X}^{\top}\mathbf{X}$ does not fit in memory. Neither regime
appeared in the tested project setting; the moderate hidden-layer sizes
accessible on a laptop ($H\le 2000$) all remained on the IRLS side in
both wall-clock time and final gap. SGPTL retains value as a sanity check (a different algorithm
converging to the same $f^{*}$ within the achievable accuracy is a
useful indication that the implementation is correct) and as a
baseline to quantify what we gain from exploiting the structure of the
LASSO problem through the MM smoothing rather than treating it as a
generic non-smooth convex programme.

\paragraph{Limitations and possible extensions.}
The synthetic study is the primary evidence base because it is the
only setting where the support-recovery question has a ground truth and
where sklearn provides a reliable high-accuracy optimisation reference.
Section~\ref{sec:real_data} validates the qualitative optimisation
behaviour on two real regression datasets (\textit{diabetes} and
\textit{California housing}) using the empirical baseline $f_{\min}$
instead: sklearn's coordinate descent stops early on diabetes and is
skipped on California. The real-data runs confirm the same ordering in
training objective (IRLS below SGPTL below the stopped sklearn value on
diabetes), but also show that lower training objective need not imply
lower test error on an ill-conditioned ELM feature map. We did
not implement an accelerated subgradient variant (for example Nesterov
smoothing of the $\ell_{1}$ term) that would close some of the rate
gap with IRLS, since the project specification fixes SGPTL as
Algorithm~A2. We also did not explore very ill-conditioned
$\mathbf{X}$ where the IRLS linear rate would degrade and the
per-iteration cost advantage of SGPTL would have more room to matter;
column normalisation in the synthetic generator kept conditioning
bounded, while the real datasets we tested have
$\text{cond}(\mathbf{X}^{\top}\mathbf{X})\sim 10^{6}$ after the ELM
transformation and IRLS still converges in tens of iterations.
